% ========================================================================
% 測圓海鏡 (Ceyuan Haijing / Sea Mirror of Circle Measurements)
% by Li Ye (李冶), 1248
% Volumes 7--9 -- EXPANDED EDITION with full Chinese mathematical text
% ========================================================================
\documentclass[11pt,a4paper]{article}

% -------- Core Packages --------
\usepackage{fontspec}
\usepackage{amsmath,amssymb,amsthm}
% Make \square work in text mode (used extensively for polynomial placeholders)
\AtBeginDocument{%
  \let\origsquare\square
  \renewcommand{\square}{\ensuremath{\origsquare}}%
}
\usepackage{geometry}
\usepackage{graphicx}
\usepackage{xcolor}
\usepackage{titlesec}
\usepackage{fancyhdr}
\usepackage{enumitem}
\usepackage{array,booktabs,longtable,multirow}
\usepackage[draft]{hyperref}
\usepackage{microtype}
\usepackage{tocloft}
\usepackage{indentfirst}
\usepackage{xeCJK}
\usepackage{etoolbox}

% -------- Page Geometry --------
\geometry{
  paperwidth=210mm,paperheight=297mm,
  left=22mm,right=22mm,top=28mm,bottom=28mm,
  headheight=14pt,footskip=30pt
}

% -------- Font Configuration --------
\setmainfont{Linux Libertine O}

% Chinese Fonts via xeCJK
\setCJKmainfont[Path=/mnt/agents/fonts/noto-sc-extracted/,BoldFont=NotoSansCJKsc-Bold.otf]{NotoSansCJKsc-Regular.otf}
\setCJKsansfont[Path=/mnt/agents/fonts/noto-sc-extracted/,BoldFont=NotoSansCJKsc-Bold.otf]{NotoSansCJKsc-Regular.otf}
\setCJKmonofont[Path=/mnt/agents/fonts/noto-sc-extracted/,BoldFont=NotoSansCJKsc-Bold.otf]{NotoSansCJKsc-Regular.otf}

\newCJKfontfamily\tradfont[Path=/mnt/agents/fonts/noto-sc-extracted/,BoldFont=NotoSansCJKsc-Bold.otf]{NotoSansCJKsc-Regular.otf}
\newCJKfontfamily\tradfontbf[Path=/mnt/agents/fonts/noto-sc-extracted/]{NotoSansCJKsc-Bold.otf}

% -------- Colors --------
\definecolor{titlecolor}{RGB}{45,55,72}
\definecolor{sectioncolor}{RGB}{55,65,81}
\definecolor{notecolor}{RGB}{120,53,15}
\definecolor{rulecolor}{RGB}{180,160,140}
\definecolor{problemcolor}{RGB}{80,40,20}
\definecolor{methodcolor}{RGB}{20,60,80}
\definecolor{workcolor}{RGB}{40,70,40}
\definecolor{anscolor}{RGB}{100,30,60}

% -------- Section Formatting --------
\titleformat{\section}
  {\normalfont\Large\bfseries\color{sectioncolor}\tradfontbf}
  {\thesection}{1em}{}
\titleformat{\subsection}
  {\normalfont\large\bfseries\color{sectioncolor}\tradfontbf}
  {\thesubsection}{1em}{}
\titleformat{\subsubsection}
  {\normalfont\normalsize\bfseries\color{sectioncolor}\tradfontbf}
  {\thesubsubsection}{1em}{}

% -------- Header/Footer --------
\pagestyle{fancy}
\fancyhf{}
\fancyhead[L]{\small\textsc{測圓海鏡}}
\fancyhead[R]{\small\thepage}
\fancyfoot[C]{\small\textit{李冶 --- 卷七至卷九}}
\renewcommand{\headrulewidth}{0.4pt}
\renewcommand{\footrulewidth}{0pt}

% -------- Custom Commands --------
\newcommand{\longline}{\noindent\rule{\textwidth}{0.4pt}}
\newcommand{\shortline}{\noindent\rule{0.5\textwidth}{0.4pt}\par}

% Problem environment
\newcommand{\owen}[1]{%
  \par\vspace{0.4em}%
  \noindent\textcolor{problemcolor}{\textbf{或問：}}{\tradfont #1}\par\nopagebreak%
}

% Method environment
\newcommand{\fayue}[1]{%
  \par\noindent\textcolor{methodcolor}{\textbf{法曰：}}{\tradfont #1}\par\nopagebreak%
}

% Working environment
\newcommand{\tsaoyue}[1]{%
  \par\noindent\textcolor{workcolor}{\textbf{草曰：}}{\tradfont #1}\par%
}

% Answer environment
\newcommand{\dayue}[1]{%
  \par\noindent\textcolor{anscolor}{\textbf{答曰：}}{\tradfont #1}\par\nopagebreak%
}

% Inline Chinese helper
\newcommand{\zh}[1]{{\tradfont #1}}

% Verse/special note
\newcommand{\shihpie}[1]{\par\noindent\textit{\tradfont #1}\par}

% Problem number marker
\newcommand{\wen}[1]{\par\vspace{0.3em}\noindent\textcolor{problemcolor}{\textbf{第#1問}}\par}

% -------- Hyperref Setup --------
\hypersetup{
  colorlinks=true,
  linkcolor=sectioncolor,
  citecolor=sectioncolor,
  urlcolor=sectioncolor,
  pdfencoding=auto,
  pdftitle={測圓海鏡 卷七至卷九（增補版）},
  pdfauthor={李冶},
  pdfsubject={Chinese Mathematics, 1248}
}

% -------- TOC Depth --------
\setcounter{tocdepth}{2}
\setcounter{secnumdepth}{2}

% -------- CJK Tweaks --------
\xeCJKsetup{PunctStyle=banjiao,CJKmath=true}
\setlength{\parindent}{2em}
\setlength{\parskip}{0.3em}

% ========================================================================
\begin{document}

% -------- Title Page --------
\begin{titlepage}
\centering
\vspace*{2cm}

{\Huge\bfseries\color{titlecolor}\tradfont 測圓海鏡}

\vspace{0.5cm}
{\Large\color{sectioncolor}\tradfont （卷七至卷九·增補全本）}

\vspace{1.5cm}
{\LARGE\tradfont 李冶\quad 撰}

\vspace{0.3cm}
{\large 1248}

\vspace{2cm}
\shortline

\vspace{1cm}
{\large\textit{Sea Mirror of Circle Measurements}}

\vspace{0.3cm}
{\normalsize Volumes 7--9 (Expanded Edition)}

\vspace{0.5cm}
{\small\textit{Transcribed and Expanded from Chinese Wikisource}}

\vspace{2cm}
{\small\tradfont
卷七：明前一十八問\quad（第105問--第122問）\\[0.3em]
卷八：明後一十六問\quad（第123問--第138問）\\[0.3em]
卷九：大斜四問、大和八問\quad（第139問--第150問）
}

\vfill
{\small Typeset with Xe\LaTeX\ using xeCJK and Noto Sans CJK SC}

\end{titlepage}

% -------- Table of Contents --------
\tableofcontents
\newpage

% ========================================================================
\section*{緒論}
\addcontentsline{toc}{section}{緒論}

\begin{tradfont}
測圓海鏡者，金元之際李治潤甫先生之所著也。先生諱治，字仁卿，號敬齋，真定欒城人，
生於金明昌元年（一一九二），卒於元至元十六年（一二七九）。此書成於淳祐八年（一二四八），
凡十二卷，計一百七十問，皆設圓城以為問，用天元術以立算，實中算之鴻篇也。

其法之要，在於天元一術。立天元一為所求之數，依題意布算，得開方式，然後開方求之。
所謂太極、天元、地元者，猶今代數之未知數也。書中凡遇開方式，皆以■記之，
蓋當時算板布列之位也。

卷七至卷九所載，凡三十四問，其題目之繁，方程式次之高，皆過於前六卷。
卷七曰「明前一十八問」，卷八曰「明後一十六問」，卷九曰「大斜四問、大和八問」。
其問題之設，或二人對行，或四處遙望，或穿城而過，或倚樹而參，錯綜變化，不可勝窮。

每問之體，例分三節：
\begin{itemize}[leftmargin=2em]
  \item \textbf{或問}：設問之辭，述人物行止之狀、所給之數。
  \item \textbf{法曰}：列方程式之術，示開方之次。
  \item \textbf{草曰}：詳布算之草，盡天元變化之妙。
\end{itemize}

又有時設\textbf{答曰}一節，直書所求之數，以為開方之驗。

書中所用十五形之名，皆有所本。通勾、通股者，大勾股形之兩邊也；
邊勾、邊股者，邊上之小形也；底勾、底股者，底下之小形也。
明勾、明股者，明朗易見之小形也；重勾、重股者，重疊相因之小形也。
至於虛勾、虛股、虛弦，則太虛之形，最為微妙。
\end{tradfont}

\vspace{0.5cm}
\noindent\textit{The \textit{Ceyuan Haijing} (測圓海鏡, ``Sea Mirror of Circle Measurements''),
completed in 1248 by the Chinese mathematician Li Ye (李冶, 1192--1279), is one
of the masterpieces of medieval Chinese algebra. The work contains 170 problems
concerning a circle inscribed in or related to a right triangle, all solved by
``tian yuan shu'' (天元術)---the method of setting up a single unknown
(``heavenly element'') to derive polynomial equations.}

\vspace{0.3cm}
\noindent\textit{Volumes 7--9 contain 34 advanced problems with higher-degree equations
(degrees 5--6), multiple unknowns, and complex geometric configurations involving
multiple observers, gate locations, and tree landmarks around the circular city.}

\newpage


% ========================================================================
\section{\tradfont 卷七\quad 明前一十八問}
\label{sec:vol7}

\begin{tradfont}
\noindent\textbf{卷七小序：}此卷一十八問，皆以明段為主。所謂明段者，
勾股形之顯然可見者也。其題多設二人對行、遙望相會之狀，
或倚門傍樹，或穿城越隍，錯綜參伍，以盡天元之變。
其方程式之高，有至三乘、四乘者，開方之法亦隨之益密。
\end{tradfont}

\vspace{0.3cm}
\begin{center}
\longline
\end{center}

\wen{一百零五}
\owen{出南門東行七十二步有樹，出東門南行三十步見之。問答同前。}
\dayue{城徑二百四十步，半徑一百二十步。}

\fayue{倍南行以乘倍東行為平實，並二行又倍之為從，一虛隅。得城徑。}

\tsaoyue{%
識別得此問名為弦外容圓，又為內率求虛積。其二行步相並為虛弦，若以相減即虛較也。
又倍東行為弦較和，倍南行即弦較較，此二數相乘則兩虛積也。若直以二行相乘，則半個虛積也。
又倍東行減於城徑，餘即二虛勾也。倍南行減於城徑則二虛股也。虛積上三事和即城徑也。

乃立天元一為圓徑，便以為三事和也。倍二行步減之，得$ \; \square \; $為黃方一，天元乘之得$ \; \square \; $為二虛積（寄左）。
然後倍東行以乘倍南行，得八千六百四十為同數，與左相消得$ \; \square \; $。
益積開平方得二百四十步，即城徑也。合問。}

\fayue{二行步相乘為實，二行步相並為從，一步虛法。得半徑。}

\tsaoyue{%
立天元一為半徑，副置二位。上加東行步得$ \; \square \; $為大差勾，下加股得$ \; \square \; $為小差股。
此二數相乘得下式$ \; \square \; $為半段黃方冪（寄左）。
然後立天元以自之，又二之，與左相消得$ \; \square \; $。
益積開平方得一百二十步，即半城徑也。}

\fayue{二雲數相乘倍之於上，加雲數差冪，權寄。並二雲數又自增乘，得數內減上位為平實，
並雲數而倍之為從，二步益隅。得半徑。}

\tsaoyue{%
立天元一為半徑，副之。上減明勾得下$ \; \square \; $為虛勾，下減股得$ \; \square \; $為虛股。
勾股相乘得$ \; \square \; $，又倍之得$ \; \square \; $，又加二行差冪$ \; \square \; $，得$ \; \square \; $為弦冪（寄左）。
然後並雲數，以自之得$ \; \square \; $於太極位，為同數，與左相消得$ \; \square \; $。
益積開平方得一百二十步，即半城徑也。}

\fayue{雲數相乘又倍之為平實，雲數相減為從，一常法。得虛勾。}

\tsaoyue{%
立天元一為虛勾。以南行減東行餘四十二步為虛較也。
以虛較加天元得$ \; \square \; $為虛股，以天元乘之得下$ \; \square \; $為直積（寄左）。
然後倍南行乘東行得$ \; \square \; $，與左相消得$ \; \square \; $。
開平方得四十八步，即虛勾也。以勾除積得九十步，即虛股也。
並勾股得$ \; \square \; $為虛和也，內加入二行並$ \; \square \; $得$ \; \square \; $，即圓徑也。}

\fayue{並二行步以自乘於上，又倍南行乘倍東行，加上位為平實，一隅法。得小和。}

\tsaoyue{%
立天元一為小和。並二行步加之得$ \; \square \; $為三事和也。
倍二行步而並之得$ \; \square \; $，以減三事和，餘$ \; \square \; $為黃方，
卻以三事和乘之，得下$ \; \square \; $為二虛積也（寄左）。
乃倍南行以乘倍東行，得$ \; \square \; $為同數，與左相消得$ \; \square \; $。
開平方得一百三十八步，即虛和也。加入二行步得二百四十步，即城徑也。合問。}

\vspace{0.5cm}
\begin{center}
\longline
\end{center}

\wen{一百零六}
\owen{丙出南門直行一百三十五步而立，甲出東門直行一十六步見之。問答同前。}
\dayue{城徑二百四十步，丙行上勾弦差八十一步。}

\fayue{%
以丙行步一百三十五再自之，得二百四十六萬零三百七十五於上。
又以甲行一十六乘丙行冪一萬八千二百二十五，得二十九萬一千六百，以乘上位，
得七千一百七十四億四千五百三十五萬為三乘方實；
以二行步相乘又倍之得四千三百二十，以乘丙行步再自之數，得一百六億二千八百八十二萬為益從；
第一廉空；以甲行乘丙行冪，得二十九萬一千六百，又倍之得五十八萬三千二百於上，
四之甲行冪一千零二十四，以乘丙行步，得一十三萬八千二百四十，減上位餘四十四萬四千九百六十為第二廉；
二行步相乘得二千一百六十為虛常法。得丙行步上勾弦差八十一。}

\tsaoyue{%
識別二數相並，得一百五十一。以減於皇極弦，餘一百三十八，即虛勾虛股並也。
若以二數相減，餘一百一十九為高弦內減平弦，又為皇極弦內少個小差弦，
又為大差弦內減個皇極弦也。

立天元一為丙行大差數。
置丙行步一百三十五，自乘得$ \; \square \; $，用天元除之，得$ \; \square \; $為勾弦並也。
上減天元得$ \; \square \; $為二丙勾也。復用丙南行乘之，得$ \; \square \; $為二積也。
又以天元除之，得$ \; \square \; $為丙勾外容圓半（泛寄）。
別置丙南行用二甲勾乘之，得$ \; \square \; $太，合用二丙勾除之。不受除，
便以此為甲股（內寄二丙勾為分母）。復用二甲勾三十二乘之，得$ \; \square \; $太為二個甲直積也。
又置丙南行內減天元，得$ \; \square \; $為黃方，以自乘得$ \; \square \; $為丙上勾弦差乘股弦差二段，
以天元除之得$ \; \square \; $為兩個丙小差也。乃用甲股乘之，得下式$ \; \square \; $。
復用丙南行除之，得$ \; \square \; $，又折半得下式$ \; \square \; $為一個甲步股弦差也，
內亦帶前二丙勾分母。復置兩個甲直積，內已寄此甲股弦差分母，便為甲步股外容圓半（寄左）。
乃再置先求到泛寄，用甲股弦差分母乘之，得$ \; \square \; $為同數，
與左相消得下式$ \; \square \; $。開三乘方得八十一步，即丙步上勾弦差也。

《鈐經》載此法，以勾弦差率冪減丙行冪，復以丙行乘之為實，以差率冪為法，如法得徑。
此法只是以勾外求圓半。合以大差除倍積，而今皆以大差冪為分母也。
依法求之，勾弦差八十一自之得六千五百六十一，以減於丙行冪一萬八千二百二十五，
餘一萬一千六百六十四，復以丙行一百三十五乘之，得一百五十七萬四千六百四十為實，
以大差冪六千五百六十一為法，如法得二百四十步，即城徑也。}

\fayue{%
二行相乘，得數又自之為三乘方實；並二行步以乘二行相乘數，又倍之為從；
二行相並數以自乘於上，又二行相減數自乘減上位為第一廉；第二廉空，一益隅。
益積開之，得半徑（其第一廉只是四段二行相乘數）。}

\tsaoyue{%
立天元一為半城徑，副置之。上加南行步得$ \; \square \; $為股，下位加東行步得$ \; \square \; $為勾。
勾股相乘得$ \; \square \; $為直積一段。以天元除之得$ \; \square \; $為弦，以自之，得$ \; \square \; $為弦冪（寄左）。
乃以勾自之，得$ \; \square \; $，又以股自之，得$ \; \square \; $，二位相並得$ \; \square \; $為同數，
與左相消，得$ \; \square \; $。益積開三乘方，得一百二十步，即半城徑也。}

\vspace{0.5cm}
\begin{center}
\longline
\end{center}

\wen{一百零七}
\owen{出東門一十六步有樹，出南門東行七十二步見之。問答同前。}
\dayue{城徑二百四十步。}

\fayue{二行步相減得數，以自之於上。又以出東門步自之，減上位為平方實，
二之出南門東行步為益從，一步常法。翻開得半徑。}

\tsaoyue{%
別得人到樹即平弦也，半圓徑即平股也。其東行七十二步則平勾平弦差也。
乃立天元一為半圓徑，加一十六減七十二，得$ \; \square \; $為勾也。
以自之得$ \; \square \; $為勾冪。又加入天元股冪得$ \; \square \; $為弦冪（寄左）。
再立天元一為半徑，加出東門步，得$ \; \square \; $即弦也。
以自之得$ \; \square \; $為同數，與左相消得$ \; \square \; $。
翻法開之得一百二十步，即半城徑也。合問。}

\vspace{0.5cm}
\begin{center}
\longline
\end{center}

\wen{一百零八}
\owen{出南門一百三十五步有樹，出東門南行三十步見之。問答同前。}
\dayue{城徑二百四十步。}

\fayue{樹去城步內減南行步，餘以為冪於上，又以樹去城步為冪內減上位為平實，
倍樹去城步為從，一虛隅。翻法得半城徑。}

\tsaoyue{%
別得人距樹即高弦也，半圓徑即高勾也，其南行三十步即高弦上小差也。
乃立天元一為半徑，加樹去城步為弦，內減小差$ \; \square \; $，得$ \; \square \; $即股也。
以自之得$ \; \square \; $為股冪，內加入天元冪，得$ \; \square \; $為弦冪（寄左）。
再置弦$ \; \square \; $以自之，得$ \; \square \; $為同數，與左相消，得式$ \; \square \; $。
翻開得一百二十步，即半城徑也。合問。}

\vspace{0.5cm}
\begin{center}
\longline
\end{center}

\wen{一百零九}
\owen{乙出東門不知遠近而立，甲出南門東行七十二步望見乙，就乙斜行一百三十六步與乙相會。問答同前。}
\dayue{城徑二百四十步，乙出東門一百二十八步。}

\fayue{以斜行步自之於上，以二行相減餘自為冪，減上位為平實，從空，一步常法。如法得半徑。}

\tsaoyue{%
別得七十二步即大差也，斜行即弦，半徑即股也。
立天元一為半徑，以自之為股冪，又以二行差六十四以自之得$ \; \square \; $為勾冪。
並二冪得$ \; \square \; $為弦冪（寄左）。然後以斜行步自之，得$ \; \square \; $為同數，
與左相消得$ \; \square \; $。開平方得一百二十步，倍之即城徑也。合問。}

\vspace{0.5cm}
\begin{center}
\longline
\end{center}

\wen{一百一十}
\owen{甲出南門不知遠近而立，乙出東門南行三十步望見甲，卻就甲斜行二百五十五步與甲相會。問答同前。}
\dayue{城徑二百四十步，甲出南門二百二十五步。}

\fayue{二行差自之為冪，以減於斜行冪為平實，一虛隅。得半徑。}

\tsaoyue{%
別得南行步即股弦差也，斜步即弦也，半徑即勾也。
乃立天元一為半城徑，以自之為冪。以二行相減餘二百二十五，以自之得$ \; \square \; $為股冪。
二冪相並得$ \; \square \; $為弦冪（寄左）。然後以斜行步自之，得$ \; \square \; $為同數，
與左相消得下$ \; \square \; $。開平方得一百二十步，即半徑也。合問。}

\vspace{0.5cm}
\begin{center}
\longline
\end{center}

\wen{一百一十一}
\owen{甲出南門東行不知步數而立，乙出東門南行三十步望見甲，斜行一百二步相會。問答同前。}
\dayue{城徑二百四十步，甲東行九十六步。}

\fayue{二行相減，餘以乘乙南行，四之於上，又加入斜行冪為平實。得虛和一百三十八。}

\tsaoyue{%
別得斜步內減南行為甲東行步也。此問以弦外容圓入之，以二行相減數乘乙南行三十步，
得$ \; \square \; $，又四之，得$ \; \square \; $為二直積也。又加入斜步冪$ \; \square \; $，
共得$ \; \square \; $即和冪也。平方而一，得一百三十八步，即虛和也。
又加斜步得二百四十步，即城徑也。合問。}

\vspace{0.5cm}
\begin{center}
\longline
\end{center}

\wen{一百一十二}
\owen{乙出東門南行不知步數而立，甲出南門東行七十二步望見乙，斜行一百二步與乙相會。問答同前。}
\dayue{城徑二百四十步，乙南行五十四步。}

\fayue{%
倍相減步以乘倍東行，得數復以減於斜步冪，餘為實。平方而一，得較也。
又以二行相減數乘倍東行為平實，以較為從方，得勾。
勾較共為長，又以斜步並入勾股共，即城徑。}

\tsaoyue{%
別得二行相減餘$ \; \square \; $為乙南行步也。以此數又減於甲東行，餘四十二步即較也。
又以二行相減數$ \; \square \; $乘倍東行得$ \; \square \; $為平實，以較為從。
平方開得四十八即勾也。勾內加較得九十步即股也。
勾股共得一百三十八，又加入斜步，共得二百四十步，即城徑也。合問。}

\vspace{0.5cm}
\begin{center}
\longline
\end{center}

\wen{一百一十三}
\owen{%
乙出南門東行，甲出東門南行，兩相望見。既而乙雲：
「我東行不及城徑一百六十八步。」甲雲：「我南行不及城徑二百一十步。」問答同前。}
\dayue{城徑二百四十步，甲南行三十步，乙東行七十二步。}

\fayue{%
半甲不及步以自之為冪，半甲不及步內減差以自之為冪。
二冪相並內卻減差冪為平實，四之甲不及內減三之乙不及，餘為益從，三步半虛法。得甲南行。}

\tsaoyue{%
別得乙不及為虛勾、半徑共，又為徑內減明勾也。
甲不及為虛股、半徑共，又為徑內減股也。
又二雲數相並為虛和、圓徑共也，雲數相減即虛較也。

乃立天元一為甲南行，以減於甲不及步又半之，得$ \; \square \; $為虛股也。
虛股內減虛較得$ \; \square \; $為虛勾。勾自之得$ \; \square \; $為勾冪也，
又股自之得下式$ \; \square \; $為股冪也。二冪相並得$ \; \square \; $為弦冪（寄左）。
然後以天元加虛較得$ \; \square \; $為乙東行。又加入天元甲南行得$ \; \square \; $為虛弦，
以自之得$ \; \square \; $為同數，與左相消得$ \; \square \; $。
開平方得三十步，即甲南行也。內加少步，即城徑也。合問。}

\vspace{0.5cm}
\begin{center}
\longline
\end{center}

\wen{一百一十四}
\owen{%
丙出南門直行，甲出東門直行，兩相望見。既而丙雲：
「我行少於城徑一百五步。」甲雲：「我行少於城徑二百二十四步。」問答同前。}
\dayue{城徑二百四十步，半徑一百二十步。}

\fayue{%
二少步相乘訖，又自乘為實；六之共步，乘雲數相乘數為益從；
十八之雲數相乘於上，又三之共步，自乘加上位，內復減丙少步冪、甲少步冪為從廉；
四十八之共步為益二廉，六十三步常法。
翻法開三乘方，得一百二十步，即半徑。}

\tsaoyue{%
別得雲數共減於倍城徑為甲丙共行數。又雲數相減即皇極差，亦為甲行不及丙行數。
立天元一為半城徑，以三之，副置二位。
上位減丙少步，得$ \; \square \; $為皇極股也，下位減甲少步得$ \; \square \; $為皇極勾也。
勾股相乘得$ \; \square \; $，以天元除之，得$ \; \square \; $為弦也。
弦自之得$ \; \square \; $為弦冪（寄左）。然後以股自之得下$ \; \square \; $為股冪於上，
又以勾自之得$ \; \square \; $為勾冪，並以加入上位，得$ \; \square \; $為同數，
與左相消得$ \; \square \; $。翻法開三乘方得一百二十步，即半城徑也。合問。}

\vspace{0.5cm}
\begin{center}
\longline
\end{center}

\wen{一百一十五}
\owen{%
甲出東門直行，乙出南門直行，各不知步數而立。
乙望見甲，就甲斜行了二百八十九步與甲相會。
其二直行共得一百五十一步。又雲甲直行少於乙直行。問答同前。}
\dayue{城徑二百四十步，甲直行一十六步，乙直行一百三十五步。}

\fayue{斜冪內減共步冪為平實，倍共步內減斜步為從，一常法。得徑。}

\tsaoyue{%
別得共數、城徑並即皇極和也。
立天元一為圓徑，加共步得$ \; \square \; $為皇極和，以自之，得$ \; \square \; $於上。
以斜行冪$ \; \square \; $減上位餘$ \; \square \; $為二直積（寄左）。
然後以天元乘斜步得$ \; \square \; $，與左相消得$ \; \square \; $。
開平方得二百四十步，即城徑也。合問。}

\vspace{0.5cm}
\begin{center}
\longline
\end{center}

\wen{一百一十六}
\owen{%
甲出東門直行，乙出東門南行，丙出南門直行，丁出南門東行，各不知步數而立。
四人遙相望，悉與城參相直。
只雲甲、丙共行了一百五十一步，乙、丁立處相距一百二步。又雲丙直行步多於甲直行步。問答同前。}
\dayue{城徑二百四十步。}

\fayue{%
共步、距步相減，得數自之於上，以共步為冪內減上為平實，
二之距步內減共步、距步差為從，一步虛法。得城徑。}

\tsaoyue{%
別得共步得城徑即皇極和也，相距步即虛弦也。
皇極和內減虛弦即皇極弦也。又共步、距步差$ \; \square \; $即皇極弦內減城徑也（此名旁差）。
乃立天元一為城徑，加共步得$ \; \square \; $為皇極和也，以自之得$ \; \square \; $於上，
以共步、距步差$ \; \square \; $加天元得$ \; \square \; $為皇極弦也，
以自之得下式$ \; \square \; $，減上位餘得$ \; \square \; $為二直積（寄左）。
然後以天元徑乘皇極弦，得$ \; \square \; $為同數，與左相消得$ \; \square \; $。
開平方得二百四十步，即城徑也。合問。}

\vspace{0.5cm}
\begin{center}
\longline
\end{center}

\wen{一百一十七}
\owen{%
甲出南門東行不知步數而立，乙出東門南行望見甲，復就甲斜行，與甲相會。
乙通計行了一百三十二步，其乙南行步不及斜行七十二步，其甲東行卻多於乙南行。問答同前。}
\dayue{城徑二百四十步，乙南行三十步，甲東行一百二步。}

\fayue{倍不及步在地，以不及步減通步以乘之為實，以四之不及步為法。得乙南行三十步。}

\tsaoyue{%
別得乙南行即股也，以減通步即虛弦也，以減不及步即虛較也，其不及步即甲東行也。
立天元一為乙南行，置不及步以天元乘之，又四之得$ \; \square \; $元為二直積（寄左）。
然後倍不及步以為弦較和於上$ \; \square \; $。
以不及步減通步得$ \; \square \; $為弦較較。
以乘上位得$ \; \square \; $太為同數，與左相消得$ \; \square \; $。
上法下實，得三十步為乙南行也。餘各以數求之。}

\fayue{%
別得通行步為兩個乙南行、一個甲東行共也。
其不及步即東行步也。雲步相並即兩個虛弦，相減即兩個乙南行也。}

\vspace{0.5cm}
\begin{center}
\longline
\end{center}

\wen{一百一十八}
\owen{%
甲出南門東行不知步數而立。乙出東門南行，望見甲，復斜行與甲相會。
二人共行了二百四步，又雲甲行不及共步一百三十二。問答同前。}
\dayue{城徑二百四十步，甲東行九十六步，乙南行一百零八步。}

\fayue{%
別得二行共即兩個虛弦也，其不及步即乙南行與一虛弦共也。
置不及步內減一弦餘三十步，即乙南行也。
以乙南行反以減虛弦，餘七十二步即甲東行也。
以乙南行減甲東行餘即虛較也。}

\tsaoyue{此問無草。}

\vspace{0.5cm}
\begin{center}
\longline
\end{center}

\wen{一百一十九}
\owen{%
乙出東門南行，甲出西門南行，甲望見乙，斜行五百一十步相會。
乙雲：「我南行少於城徑二百一十步。」問答同前。}
\dayue{城徑二百四十步，乙南行三十步。}

\fayue{少步冪為平實，四斜步內減二少步為益從，五步常法。得乙南行。}

\tsaoyue{%
別得少步為徑內減股。
立天元一為乙南行，以二之減於倍斜行步，得$ \; \square \; $為梯底也。
以二之天元乘之，得$ \; \square \; $為徑冪（寄左）。
再置天元加少步，得下式$ \; \square \; $為城徑，以自之得$ \; \square \; $，
與左相消得$ \; \square \; $。開平方得三十步，即乙南行也。加少步即城徑也。合問。}

\vspace{0.5cm}
\begin{center}
\longline
\end{center}

\wen{一百二十}
\owen{%
乙出南門東行，甲出北門東行，甲望見乙，斜行二百七十二步與乙相會。
乙雲：「我東行不及城徑一百六十八步。」問答同前。}
\dayue{城徑二百四十步，乙東行七十二步。}

\fayue{以不及步冪之為實，四斜內減二之不及步為虛從，五常法。平開得乙東行七十二步。}

\tsaoyue{%
別得不及步為城徑減明勾也。
立天元一為乙東行，以倍之減於二之斜行步，得下$ \; \square \; $為梯底也。
倍天元乘之，得$ \; \square \; $為徑冪（寄左）。
再置天元加不及步，得$ \; \square \; $為城徑，以自之得$ \; \square \; $為同數，
與左相消得$ \; \square \; $。開平方得七十二步，即乙東行也。加入少步即城徑也。合問。}

\vspace{0.5cm}
\begin{center}
\longline
\end{center}

\wen{一百二十一}
\owen{%
乙出南門東行，丁出東門南行，卻有甲丙二人共在西北隅，甲向東行，丙向南行，
四人遙相望見，俱與城參相直。既而相會，甲雲：「我多乙二百四十八步。」
丙雲：「我多於丁五百七十步。」問答同前。}
\dayue{城徑二百四十步。}

\fayue{二多步相乘為平實，並二多步而半之為從，七分半常法。得城徑。}

\tsaoyue{%
別得甲多步為大勾內減明勾也，丙多步為大股內少股也。
又乙東行得一虛勾為半徑，丁南行得一虛股為半徑。
又二多數相並得$ \; \square \; $為大和內少虛弦也，又二少數相減餘$ \; \square \; $為兩個角差。
又甲多步內減半徑即勾方差也，丙多步內減半徑即股方差也。

立天元一為城徑，以半之減於甲多步得$ \; \square \; $為勾方差，
又以半徑減於丙多步得$ \; \square \; $為股方差。
二差相乘得$ \; \square \; $為徑冪（寄左）。然後以天元冪與左相消，得下式$ \; \square \; $。
開平方得二百四十步，即城徑也。合問。}

\vspace{0.5cm}
\begin{center}
\longline
\end{center}

\wen{一百二十二}
\owen{%
甲丙二人俱在西北隅，甲向東行，丙向南行。
又乙出南門東行，丁出東門南行，各不知步數而立。
四人遙相望見，悉與城參相直。既而相會，甲雲：
「我與乙共行了三百九十二步。」丙雲：「我與丁共行了六百三十步。」問答同前。}
\dayue{城徑二百四十步。}

\fayue{%
甲乙共自之為冪，丙丁共自之為冪，二冪又相乘為三乘方實。
甲乙共自之為冪，以丙丁共乘之於上。
又以丙丁共自之為冪，以甲乙共乘之加上位為益從。
甲乙共自之為冪，丙丁共自之為冪，並，以七分半乘之於上。
又以甲乙共乘丙丁共，得數減上位為第一益廉。
並二共數，以七分半乘之為第二廉。
以七分半自之，得五分六厘二毫五絲於上位。
以一步內減上位，餘四分三厘七毫五絲為虛隅。得城徑。}

\tsaoyue{%
別得甲為大勾，乙為明勾，丙為大股，丁為股也。
甲乙共內減半徑即是黃長弦也，丙丁共內減半徑即黃廣弦也。
黃長弦、黃廣弦二數相減，餘為兩個皇極差也。

乃立天元為城徑，半之副置二位。
上以減於甲乙共數，得$ \; \square \; $即黃長弦也，以自之得$ \; \square \; $為黃長弦冪也，
內減天元一冪，餘得下式$ \; \square \; $為勾方差冪也。
下位以減於丙丁共，得下式$ \; \square \; $即黃廣弦也，
以自之得$ \; \square \; $為黃廣弦冪也，內減天元一冪，餘得$ \; \square \; $為股方差冪也。
再以勾方差冪、股方差冪相乘，得$ \; \square \; $為徑冪（寄左）。
然後以天元為冪，又以冪自之，與左相消得下式$ \; \square \; $。
開三乘方得二百四十步，即城徑也。合問。}

\newpage


% ========================================================================
\section{\tradfont 卷八\quad 明後一十六問}
\label{sec:vol8}

\begin{tradfont}
\noindent\textbf{卷八小序：}此卷一十六問，承卷七之緒，仍以明段為主，
而益之以極段、虛段之變。所謂極段者，皇極勾股形之各邊也；
虛段者，太虛勾股形之各邊也。此三形交錯互見，開方之術益臻精妙。
其方程式之高，有至四乘方、五乘方者，實全書之樞紐也。
\end{tradfont}

\vspace{0.3cm}
\begin{center}
\longline
\end{center}

\wen{一百二十三}
\owen{%
出南門向東有槐樹一株，出東門向南有柳樹一株。
丙丁俱出南門，丙直行，丁往至槐樹下。
甲乙俱出東門，甲直行，乙往至柳樹下。
四人遙相望見，各不知所行步數。
只雲丙丁共行了二百七步，甲乙共行了四十六步，
又雲甲丙立處相距二百八十九步。問答同前。}
\dayue{城徑二百四十步，半徑一百二十步。}

\fayue{以二共相減數又以減距數為實，二為法。得平勾。}

\tsaoyue{%
識別得丙丁共即明和也，甲乙共即和也，相距步即極弦也。
二共相並即極弦內少個虛黃也，又為極和內少個虛和也。
二共相減餘為平勾高股差也，又為虛差極差共也，又為通差內減極差也。

立天元為平勾，加入二共相減數，得$ \; \square \; $為高股，又加天元得$ \; \square \; $為極弦（寄左）。
以相距步二百八十九與左相消，得$ \; \square \; $。
上法下實，如法得六十四，即平勾也。
以二共相減數加平勾得二百二十五為高股，復以平勾乘之，得一萬四千四百步。
開平方得一百二十步，即城半徑也。合問。}

\fayue{二共數並，以減相距數，餘者半之為泛率。以泛率加丙丁共為長，以泛率加甲乙共為闊。
長闊相乘為平方實。得半徑。}

\tsaoyue{%
置極弦內減二共並數，餘三十六步即虛黃也。
半之，副置二位。上以加明和得二百二十五步為高股也，下以加和得六十四步為平勾也。
二位相乘得一萬四千四百步。開平方得一百二十步，即半徑也。合問。}

\vspace{0.5cm}
\begin{center}
\longline
\end{center}

\wen{一百二十四}
\owen{%
依前見丙丁共二百七步，甲乙共四十六步。又雲二樹相去一百二步。問答同前。}
\dayue{城徑二百四十步。}

\fayue{%
以甲乙共乘樹相去步，得數又以自之為平實；從空；
並二共數為冪於上，內減甲乙共自之數、丙丁共自之數為益隅。得弦。}

\tsaoyue{%
識別得兩樹相去步即虛弦也。餘數具前。
立天元一為弦。置明和以天元乘之，合和除，不除，便以$ \; \square \; $元為明弦也（內帶和分母）。
乃置虛弦以分母和乘之，得$ \; \square \; $太，加入明弦，得$ \; \square \; $為極股也，內帶和分母。
以自之得下式$ \; \square \; $為極股冪（內寄和冪為分母）。
又以天元加虛弦得下$ \; \square \; $為極勾，以自之得$ \; \square \; $。
又以和冪$ \; \square \; $乘之，得$ \; \square \; $為勾冪也。
勾股冪相並得$ \; \square \; $為兩積一較冪也，內有和冪分母（寄左）。
然後置明弦$ \; \square \; $元於上，以和乘天元加上位，得$ \; \square \; $元為二弦並也。
又置虛弦以和乘之得$ \; \square \; $太，並入上位，得下式$ \; \square \; $為極弦，
以自之得$ \; \square \; $為同數，與左相消得$ \; \square \; $。
開平方得三十四步，即弦也。}

\fayue{以樹相去步自之，又以甲乙共乘之為平實，從空，倍丙丁共為虛隅。得弦。}

\tsaoyue{%
立天元一為弦。依前術求得明弦$ \; \square \; $元，便以為皇極勾弦差也（內帶和分母）。
以天元弦便為皇極股弦差，以乘之，又倍之，得$ \; \square \; $為虛弦冪（內有和分母。寄左）。
然後以虛弦自之，又以分母$ \; \square \; $乘之，得四十七萬八千五百八十四為同數，
與左相消得$ \; \square \; $。開平方得三十四步，即弦也。合問。}

\vspace{0.5cm}
\begin{center}
\longline
\end{center}

\wen{一百二十五}
\owen{皇極大小差共一百八十七步，明黃、黃共六十六步。問答同前。}
\dayue{太虛黃方三十六步，城徑二百四十步。}

\fayue{後數自乘為實，前後數相減，餘為法。得虛黃方三十六。}

\tsaoyue{%
別得一百八十七即明弦二弦共也，其六十六即太虛大小差共也。
又二數相並，得$ \; \square \; $即明、二和共。若以相減，餘$ \; \square \; $即明、四差共也。

立天元一為太虛黃方面，加二黃共，得$ \; \square \; $即虛弦也。
倍虛弦又加天元，得$ \; \square \; $即城徑也。
又以虛弦加皇極大小差，得$ \; \square \; $即極弦也，
以極弦乘城徑得$ \; \square \; $，為兩段皇極勾股積（寄左）。
再以極弦虛弦相並，得$ \; \square \; $，即皇極勾股共也，
自之得$ \; \square \; $，內減皇極弦冪$ \; \square \; $，得$ \; \square \; $為同數，
與寄左相消得$ \; \square \; $。上法下實，如法得三十六步，即太虛黃方面也。合問。}

\vspace{0.5cm}
\begin{center}
\longline
\end{center}

\wen{一百二十六}
\owen{%
東門南有柳一株，南門東有槐一株。
甲出東門直行，丙出南門直行。
甲、丙、柳、槐悉與城參相直，
既而甲就柳樹斜行三十四步至柳樹下，丙就槐樹斜行一百五十三步至槐樹下。問答同前。}
\dayue{城徑二百四十步，虛弦一百二步。}

\fayue{%
雲數相乘，倍之便為平方實，開方得虛弦一百二步。
以此弦加甲行步即極勾，以此弦加丙行步即極股。餘各依法求之。}

\tsaoyue{%
識別：甲斜行即弦也，丙斜行即明弦也。
立天元一為虛弦。置甲乙共以天元減之，得$ \; \square \; $為虛和也。
以天元減甲乙共餘$ \; \square \; $即虛勾虛股共也。
別置虛弦加甲斜行得$ \; \square \; $即極勾，虛弦加丙斜行得$ \; \square \; $即極股也。
極勾極股相乘得$ \; \square \; $，以天元虛弦除之得$ \; \square \; $為極弦也。
以自增乘得$ \; \square \; $為極弦冪（寄左）。
乃以極勾自之得$ \; \square \; $為勾冪，極股自之得$ \; \square \; $為股冪，
二冪相並得$ \; \square \; $為同數，與左相消得$ \; \square \; $。
開平方得一百二步，即虛弦也。合問。}

\vspace{0.5cm}
\begin{center}
\longline
\end{center}

\wen{一百二十七}
\owen{%
東門南有柳一株，南門東有槐一株。
甲出東門直行，丙出南門直行，二人遙相望，槐柳與城邊悉相直。
既而甲復斜行至柳樹下，丙復斜行至槐樹下，各不知步數。
只雲丙共行了二百八十八步，甲斜行與柳至東門步共得六十四步。問答同前。}
\dayue{甲直行一十六步，城徑二百四十步。}

\fayue{%
二雲數相乘於上，以六十四步自之，又二之，減上位為平實。
四之六十四於上，倍丙行內減上位為從。二十常法。得甲直行步一十六。}

\tsaoyue{%
別得丙共步即明股、明弦和也，六十四即平勾也，內甲斜行即弦也，柳至東門步即股也。
又二雲數相並即明差與極弦共也，二雲數相減即明差與平勾高股差共也。
又平勾內減勾即虛勾也。

立天元一為勾。置丙共步以天元乘之，復以六十四除之得$ \; \square \; $為明勾也。
又以天元減於六十四，得$ \; \square \; $為虛勾也，並虛明二勾$ \; \square \; $為半徑也。
以自之得$ \; \square \; $，倍之得$ \; \square \; $為半段圓城徑冪（寄左）。
乃以天元加六十四，得$ \; \square \; $為勾圓差於上；
又以明勾加丙共步，得$ \; \square \; $為股圓差於下。
上下相乘得$ \; \square \; $為同數，與左相消得$ \; \square \; $。
開平方得一十六步，即勾也。此勾乃甲出東門直行步也。餘皆依數求之。合問。}

\vspace{0.5cm}
\begin{center}
\longline
\end{center}

\wen{一百二十八}
\owen{%
東門南有柳樹一株，南門東有槐樹一株。
甲出東門直行，丙出南門直行，二人遙相望，槐柳與城邊悉相直。
既而甲復斜行至柳樹下，丙復斜行至槐樹下，各不知步數。
只雲甲共行五十步，丙斜行與槐至南門步共得二百二十五步。問答同前。}
\dayue{城徑二百四十步，丙直行一百三十五步。}

\fayue{%
以二百二十五步自之為冪，又以此冪自為冪於上。
置甲共行以二百二十五步三度乘之，得數復折半，減上位為平實。
置二百二十五步自之數，以二雲數相減數乘之，又倍之於上。
倍五十步在地，以二百二十五步自之數乘之，復折半加上位為益從。
雲數相減自乘於上，以雲數相乘，復折半，減上位為常法。得明股。}

\tsaoyue{%
識別得甲共步即勾、弦共也，二百二十五即高股也，內丙斜行即明弦，槐至南門步即明勾也。
又二雲數相並即極弦內減一個差也，雲數相減即差與高股、平勾差共也。
又高股內減明股即虛股也。

立天元一為明股，即丙出南門直行步也。
置五十步以天元乘之得$ \; \square \; $元，合高股除。
不除，便以此$ \; \square \; $元為股也，內帶高股$ \; \square \; $分母。
再置高股內減天元得$ \; \square \; $為虛股，以分母高股乘之，得下式$ \; \square \; $，
加入股得$ \; \square \; $即半徑也。
以自增乘得下$ \; \square \; $為半徑冪也。內帶高股冪為母（寄左）。
然後置甲共步以分母高股乘之，得$ \; \square \; $太，加入股得$ \; \square \; $為勾圓差於上（內帶高股分母）。
又以天元加高股得$ \; \square \; $為股圓差於下。
上、下相乘得$ \; \square \; $。又以分母高股乘之，得$ \; \square \; $，
復折半得$ \; \square \; $為同數，與左相消得$ \; \square \; $。
開平方得一百三十五步，即明股也。合問。}

\vspace{0.5cm}
\begin{center}
\longline
\end{center}

\wen{一百二十九}
\owen{通勾、通弦共一千步，勾、弦共五十步。問答同前。}
\dayue{城徑二百四十步，股三十步。}

\fayue{%
置一千減二之五十步為泛率。以自乘，復半之於上。
又置泛率復以五十乘之，加上位為平實。
二十二之泛率於上。以四十二乘五十，得數內減泛率，加上位為益從。
二百為常法。得股。}

\tsaoyue{%
立天元一為股。
置一千以天元乘之，以五十除之，得$ \; \square \; $元為通股也。
又以天元加五十步，得$ \; \square \; $即小差也，通股加小差得$ \; \square \; $即通弦也。
以通弦減一千得$ \; \square \; $，即通勾也。
以小差減通勾得$ \; \square \; $，即圓徑也。
以圓徑減通股得$ \; \square \; $即大差也。
置大差以小差乘之得$ \; \square \; $（寄左）。
然後置圓徑以自之得$ \; \square \; $，折半得$ \; \square \; $，與左相消得$ \; \square \; $。
開平方得三十步，即股也。合問。}

\vspace{0.5cm}
\begin{center}
\longline
\end{center}

\wen{一百三十}
\owen{通勾、通弦共一千步，明勾、明弦共二百二十五步。問答同前。}
\dayue{城徑二百四十步，明股一百三十五步。}

\fayue{%
以後數再自乘，又以前數乘之為平實；以後數為冪，復以前數乘之為從；
以前數冪為虛常法。得明股。}

\tsaoyue{%
別得二百二十五步即高股也。
立天元一為明股。置一千以天元乘之，合以高股除不受除，
便以此一零零零元為通股（內帶高股為母）。
以天元加高股，得$ \; \square \; $即大差也。
置大差以高股分母乘之，得$ \; \square \; $，即帶分大差也。
以此減於通股，餘$ \; \square \; $即圓徑也。
以自增乘，得$ \; \square \; $（寄左。內帶高股冪分母）。
然後置一千以高股分母通之，得$ \; \square \; $太，內減帶分大差，
得$ \; \square \; $為兩個通勾也。
內減兩個圓徑得$ \; \square \; $，為兩個小差也。
以帶分大差乘之，得下式$ \; \square \; $為同數，與左相消得$ \; \square \; $。
開平方得一百三十五步，即明股也。合問。}

\vspace{0.5cm}
\begin{center}
\longline
\end{center}

\wen{一百三十一}
\owen{通股、通弦共一千二百八十步，股、弦共六十四步。問答同前。}
\dayue{城徑二百四十步，勾一十六步。}

\fayue{%
雲數相乘為平實，前數為益從。
置前數以後數除之，得二十為泛率。
泛率減一以自乘於上，又倍泛率減一加上位為常法。倒積開得勾一十六。}

\tsaoyue{%
別得六十四步即平勾也。
立天元一為勾。置前數以天元乘之，以後數除之，得$ \; \square \; $元即通勾也。
又置天元加後數，得$ \; \square \; $即小差也。
以小差減通勾，餘$ \; \square \; $即圓徑也。
以自之得$ \; \square \; $（寄左）。
然後以小差減於前數，得$ \; \square \; $為二通股。
內減兩個圓徑，得$ \; \square \; $為二大差也。
以小差乘之得下$ \; \square \; $，與左相消得$ \; \square \; $。
開平方得一十六步，即勾也。合問。}

\vspace{0.5cm}
\begin{center}
\longline
\end{center}

\wen{一百三十二}
\owen{通股、通弦共一千二百八十步，明股、明弦共二百八十八步。問答同前。}
\dayue{城徑二百四十步，明勾七十二步。}

\fayue{%
二數相減，以後數乘之，內減後數冪，又半之為泛率。以自之為平實。
置前數加二之後數而半之為次率，以乘泛率，倍之於上，
以後數乘泛率減上位為益從。
次率自乘於上，以前數加次率，復以後數乘之，減上位為隅法。得明勾。}

\tsaoyue{%
別得二數相減，餘$ \; \square \; $為通勾、通股及明勾共也。
立天元一為明勾。置前數以天元乘之，合以後數除之。
不除，便以此$ \; \square \; $元為通勾也（內寄後數分母）。
又以二數相減，得數內又減天元得$ \; \square \; $為通和也。
乃以分母二百八十八之，得下式$ \; \square \; $，內減通勾，餘$ \; \square \; $為通股也。
又以天元加後數，又以分母（即後數也）通之，得$ \; \square \; $為大差也。
以此大差減於通股，得下式$ \; \square \; $，為一個圓徑也。
半之得$ \; \square \; $，以自之得$ \; \square \; $為半徑冪（寄左）。
然後以半圓徑減通勾，得$ \; \square \; $為底勾，又以天元乘之，
又以分母二百八十八之，得$ \; \square \; $為同數，與左相消得$ \; \square \; $。
開平方得七十二步，即明勾也。合問。}

\vspace{0.5cm}
\begin{center}
\longline
\end{center}

\wen{一百三十三}
\owen{%
明股、明弦並二百八十八步，勾、弦並五十步。
又雲明股、勾並多於虛弦四十九步。問答同前。}
\dayue{城徑二百四十步。}

\fayue{前二數相並，內減二之多步，即圓徑。又只以前二數相乘，便是半徑冪。}

\tsaoyue{%
識別得前二數相減而半之，即極差也。其多步名傍差，又為圓徑不及極弦數。
立天元一為半徑。置前二數相乘得$ \; \square \; $，以天元除之得$ \; \square \; $為極弦也。
以天元加多步得$ \; \square \; $為虛弦。極弦內減虛弦餘$ \; \square \; $為旁差，與多步相應。
又以極弦自之得$ \; \square \; $，內減虛弦冪$ \; \square \; $，餘$ \; \square \; $為二積。
以天元除之得$ \; \square \; $為勾股和，即前二數之並也，以此為同數。
與左相消得$ \; \square \; $，開平方得一百二十步，即半徑也。
倍之得二百四十步，即城徑也。合問。}

\vspace{0.5cm}
\begin{center}
\longline
\end{center}

\wen{一百三十四}
\owen{平差、高差共一百六十一步，明股、勾並多於虛弦四十九步。問答同前。}
\dayue{城徑二百四十步，平勾六十四步。}

\fayue{二數相減又半之，以自乘為實，後數為法。得平勾。}

\tsaoyue{%
立天元一為平勾。以加前數得$ \; \square \; $為高股也。
又以天元加高股得$ \; \square \; $為極弦，內減後數得$ \; \square \; $，
又半之，得$ \; \square \; $為半徑。
以自之，得$ \; \square \; $（寄左）。
然後以天元乘高股，得$ \; \square \; $為同數，與左相消得$ \; \square \; $。
上法下實，得六十四步，即平勾也。合問。}

\vspace{0.5cm}
\begin{center}
\longline
\end{center}

\wen{一百三十五}
\owen{%
平勾、高股差一百六十一步，明差、差並七十七步。
又雲極弦多於城徑四十九步。問答同前。}
\dayue{城徑二百四十步，半徑一百二十步。}

\fayue{%
並上二位而半之為平率。其四十九即旁率也。
副置平率，上加旁率，下減旁率，以相乘為實。倍旁差為法。得勾圓差。}

\fayue{%
又法：並上二位而半之，為平率。其四十九即旁率也。
副置旁率，上以減於平率，下以減於前數，以相乘為實。倍旁差為法。得勾圓差。}

\fayue{%
又法：求半徑：副置平率，上加旁率，下減旁率，以相乘為實，倍旁差為法。得半徑。}

\tsaoyue{%
立天元一為半徑，又為半之股圓差上弦較較，又為半之勾圓差上弦較和也。
內減勾圓差上勾股較$ \; \square \; $，餘$ \; \square \; $為半之勾圓差上弦較較也。
置股圓差上勾股較$ \; \square \; $，以半之勾圓差上弦較較乘之，得$ \; \square \; $（寄左）。
然後以半之股圓差上弦較較乘勾圓差上勾股較，
得$ \; \square \; $元為同數，與左相消得下式$ \; \square \; $。
上法下實，得一百二十步，即半徑也。合問。}

\tsaoyue{%
識別得平勾、高股差名角差。
副置角差，上加七十七而半之，得$ \; \square \; $，即極差也。
下減七十七而半之，得$ \; \square \; $，即虛差也。
角差加極差得$ \; \square \; $，即通差也。
又極弦多於城徑步名為旁差。
副置角差，上加旁差得$ \; \square \; $，為兩個高段上勾股較，
下減旁差得$ \; \square \; $為兩個平段上勾股較也。
又副置極差，上加旁差得$ \; \square \; $為股圓差上勾股較，
下減旁差$ \; \square \; $為勾圓差上勾股較也。
立天元一為勾圓差。依法求得通差，加入天元得$ \; \square \; $即大差也。
以天元乘之得$ \; \square \; $為半段圓徑冪（寄左）。
乃置大差$ \; \square \; $內減股圓差上勾股較$ \; \square \; $，
餘有$ \; \square \; $為股圓差之勾於上。
再置天元內加勾圓差上勾股較$ \; \square \; $，得$ \; \square \; $為勾圓差之股。
以乘上位得$ \; \square \; $為同數，與左相消得$ \; \square \; $。
上法下實，得八十步，即勾圓差也。}

\vspace{0.5cm}
\begin{center}
\longline
\end{center}

\wen{一百三十六}
\owen{%
又依前問：見角差一百六十一步，見明差差並七十七步，又見太虛弦較較六十步。問答同前。}
\dayue{城徑二百四十步，平勾六十四步。}

\fayue{%
前二數相減而半之，得數加入半之太虛弦較較為泛率，以自乘為平實。
置一百六十一內減二之泛率為從，一常法。得平勾。}

\tsaoyue{%
別得$ \; \square \; $即二股也。
立天元一為平勾。先以前二數相減而半之，得$ \; \square \; $為虛差。
以虛差加股得$ \; \square \; $，即明勾也。
以明勾加天元得$ \; \square \; $，為平弦，以自之得$ \; \square \; $，
內減天元冪，得$ \; \square \; $為半徑冪（寄左）。
然後以天元加一百六十一為高股，以天元乘之，得$ \; \square \; $為同數，
與左相消得$ \; \square \; $。開平方得六十四步，即平勾也。}

\fayue{%
又法曰：前數內加半之太虛弦較較，以自乘，內減前數自乘為實，
前數內減太虛弦較較為從，一常法。
開平方得平勾六十四。此更不用明差差並也。}

\tsaoyue{%
依前求平勾。前高股內加股，得$ \; \square \; $為高弦也。
以自之得$ \; \square \; $於上位，內減高股冪$ \; \square \; $，餘得$ \; \square \; $為半徑冪（寄左）。
然後以天元乘高股，得$ \; \square \; $為同數，與左相消得下$ \; \square \; $。
開平方得六十四步，即平勾也。合問。}

\vspace{0.5cm}
\begin{center}
\longline
\end{center}

\wen{一百三十七}
\owen{高差、平差並一百六十一步，明差、差並七十七步。問答同前。}
\dayue{城徑二百四十步。}

\fayue{%
以前數自乘於上，二數相並而半之，以自乘減上位，得數復自增乘為平實。
前數自之於上，又以四之前數乘之，寄位。
以前數自之於上，並二數而半之，以自乘減上位，得數又以四之前數乘，又倍之，減於寄位為從。
前數自之，又四之於上。又以四之前數為冪，加上位，權寄。
以前數為冪於上，並二數而半之，以自乘減上位，得數復八之於上。
又以四之前數為冪加入上位，並以減於權寄為常法。得平勾。}

\tsaoyue{%
識別得二位相並而半之，得$ \; \square \; $，即極差也。
立天元一為平勾，加一百六十一得$ \; \square \; $為高股，高股內又加天元，得$ \; \square \; $為極弦。
以自之得$ \; \square \; $於上，內減極差冪一萬四千一百六十一，
餘$ \; \square \; $為兩段極積。合以極弦除，不除寄為母，便以此為城徑。
以自增乘得$ \; \square \; $為圓徑冪（內有極弦冪分母。寄左）。
然後以天元乘高股，又四之得$ \; \square \; $，又以分母極弦冪$ \; \square \; $通之，
得$ \; \square \; $為同數，與左相消得$ \; \square \; $。
開平方得六十四步，即平勾也。合問。}

\vspace{0.5cm}
\begin{center}
\longline
\end{center}

\wen{一百三十八}
\owen{見明和二百七步，和四十六步。問答同前。}
\dayue{城徑二百四十步，小差四步。}

\fayue{%
二和上下相減，數同則止，名為泛率。
又以二和直相減，餘為泛實（此則角差也）。
乃以泛率除泛實，所得為差率也。
以差率加減泛率，若半訖，與勾股相應者，其泛率便為和率，
其泛實便為較率乘和率也。
若不相應，則直取差率以消息之，定為相管和率
（其勾股數少，得見弦黃而相為率者。
勾三股四，則其和七，而其較一也。
勾五股十二，則其和一十七，而其較七也。
勾八股十五，則其和二十三，而其較亦得七也。
勾七股二十四，則其和三十一，而其較一十七也。
勾九股四十，則其和四十九，而其較三十一也。
此消息之大略也。餘皆仿此）。
乃以和率約二和，其明和所得為明壘率，其和所得為壘率也。
又副置和率，上加差率而半之，則為股率也；
下位減差率而半之，則為勾率也。
既見勾、股及差三率，各以壘率乘之，即各得勾、股及差之真數也。}

\fayue{%
又法：二雲數相並，以自乘於上。
二雲數相乘，又四之以減上位為實。
二雲數相並，以六步半乘之於上，又二數相並，以四步半乘之，又四之，
以並入上位為從方。
以七十步零四分三厘七毫五絲為常法。得小差四步。}

\tsaoyue{%
以二和相約分得率一、明率四步半，其兩數大小差率並同。
又別得明小差、大差俱為半虛黃也。
立天元一為小差，以四步半乘之得$ \; \square \; $為大差也，又為明小差，又為半虛黃。
置此大差又以四步半乘之得$ \; \square \; $為明大差也。
其四差相並得$ \; \square \; $，減於二和並，得$ \; \square \; $即兩段太虛大小差並也。
內加三段虛黃方$ \; \square \; $，得$ \; \square \; $，合成一個太虛三事和，即圓城徑也。
以自增乘得$ \; \square \; $為徑冪（寄左）。
乃置和加半虛黃，得$ \; \square \; $為平勾。
又置明和內加半虛黃，得$ \; \square \; $為高股，
勾股相乘得下式$ \; \square \; $，又四之得$ \; \square \; $為同數，
與左相消得下式$ \; \square \; $。開平方得四步，即小差也。合問。}

\vspace{0.5cm}
\begin{center}
\longline
\end{center}

\wen{一百三十九}
\owen{明勾二股共八十八步，明股二勾共一百六十五步。問答同前。}
\dayue{城徑二百四十步，半虛黃方一十八步。}

\fayue{%
先識別得二大差共、二小差共及四差共。
乃以二大差、二小差相乘為實，以四差共為法。如法得半之虛黃方一十八。}

\tsaoyue{%
先置前後雲數以約法約之，得一十一即壘率也。
復各置前後數如壘率而一，前得八即勾率也，後得一十五即股率也。
再以勾、股率求得較率七，和率二十三，弦率一十七，黃方率六，大差率九，小差率二。
既見諸率，各以壘率乘之，其二和共得$ \; \square \; $，二較共得$ \; \square \; $，
二弦共得$ \; \square \; $，二黃共得$ \; \square \; $，二大差共$ \; \square \; $，
二小差共$ \; \square \; $，四差共$ \; \square \; $。
已上皆為明所得之共數也。

乃立天元一為半虛黃，便為明小差，又為大差也。
以減於大差共，得$ \; \square \; $即明大差也。
又以減於小差共，得$ \; \square \; $即小差也。
以二數相增乘，得$ \; \square \; $（寄左）。
以天元冪與寄左相消，得$ \; \square \; $。
上法下實，得一十八步，即半之虛黃方也。
以倍之得$ \; \square \; $，又加於二黃共六十六共得一百二，
即明勾、股共也，又為極黃方，又為虛弦也。
又以三十六減於一百八十七，餘一百五十一，即明股勾共也。
此數內減虛弦，餘$ \; \square \; $為明二差較也，此名旁差。
以旁差減二弦共一百八十七，餘得$ \; \square \; $，即太虛和也。
卻加入虛弦一百二，並得$ \; \square \; $，為太虛三事和，即圓城徑也。合問。}

\fayue{%
又法：以虛黃方加於二和共二百五十三，得$ \; \square \; $為極弦也。
以旁差減極弦餘二百四十步。亦同。}

\tsaoyue{%
前後副置勾、股、較、和、弦、黃六率在地，
前以小差率二因之，則勾得$ \; \square \; $，股得$ \; \square \; $，
較得$ \; \square \; $，和得$ \; \square \; $，弦得$ \; \square \; $，黃得$ \; \square \; $，即段各數也。
後以大差率九因之，則勾得$ \; \square \; $，股得$ \; \square \; $，
較得$ \; \square \; $，和得$ \; \square \; $，弦得$ \; \square \; $，黃得$ \; \square \; $，即明段各數也。
既得明、各數，餘皆可知。}

\newpage


% ========================================================================
\section{\tradfont 卷九\quad 大斜四問、大和八問}
\label{sec:vol9}

\begin{tradfont}
\noindent\textbf{卷九小序：}此卷凡十二問，分為二節。
上節曰「大斜四問」，以通勾股形之斜邊為主，方程式之高有至五乘、六乘者，
實全書開方之極致。下節曰「大和八問」，以通勾股形之三事和為主，
其術益奧，其變益奇。學者能精於此，則天元之術思過半矣。
\end{tradfont}

\subsection{\tradfont 細草卷九上\quad 大斜四問}

\begin{center}
\longline
\end{center}

\wen{一百四十}
\owen{%
甲丙俱在中心，丙望南門直行，不知步數而止。
甲出東門直行，不知步數望見丙，斜行與丙相會。
二人共行了六百八十步，仍雲甲直行少於丙直行一百一十九步。問答同前。}
\dayue{城徑二百四十步，皇極勾一百三十六步。}

\fayue{%
二數相減，餘以為冪，內卻減差冪為平實；
二數相減又四之於上，又加入二之差步為益從；二步常法。得皇勾一百三十六。}

\tsaoyue{%
別得共步即皇極三事和，少步即勾股差也。
立天元一為皇極勾。加少步得$ \; \square \; $為股也，又以天元加股得$ \; \square \; $為和也。
以和減共步得$ \; \square \; $為弦也，弦自之得$ \; \square \; $為一段弦冪（寄左）。
然後置股以天元乘之，又倍之，得$ \; \square \; $為二直積，
如入少步冪$ \; \square \; $共得$ \; \square \; $為同數，與左相消得$ \; \square \; $。
平方而一，得一百三十六即勾也。勾加差為股，勾股相乘，倍之為實，
勾股和減共步為法。得城徑。}

\fayue{%
又法：和數與倍差相加、相減，二得數相乘為平實；
雲數並與雲數差相並得數，以減於八之共步為益從；一步常法。得皇極黃方一百二。}

\tsaoyue{%
立天元一為黃方（即虛弦也），副置之。
上位加共步，得$ \; \square \; $為二和也；下位減共步，得$ \; \square \; $為二弦也。
先以二和自乘，得$ \; \square \; $為四段和冪。
又以二弦自乘，得$ \; \square \; $為四段弦冪。
二數相減，餘得$ \; \square \; $元，又倍之得下式$ \; \square \; $元，
為十六段直積於天元位（寄左）。
然後副置二和，上位加二之少步，得$ \; \square \; $為四股。
下位減二之少步，得$ \; \square \; $為四勾。
勾股相乘，得$ \; \square \; $為同數，與左相消得$ \; \square \; $。
平方而一，得一百二步，即皇極黃方也。餘各依法求之。合問。}

\vspace{0.5cm}
\begin{center}
\longline
\end{center}

\wen{一百四十一}
\owen{%
甲丙俱在西北隅起，丙向南行不知步數而立。
甲向東行望見丙，就丙斜行六百八十步與丙相會。
丙雲：「我南行步多於甲東行二百八十步。」問答同前。}
\dayue{城徑二百四十步，大勾三百二十步。}

\fayue{以雲數差乘雲數並為實，倍多步為從，二為平隅。得大勾三百二十。}

\tsaoyue{%
立天元為大勾。加差得$ \; \square \; $為股，倍天元乘之，得$ \; \square \; $為二積（寄左）。
然後以斜步、多步並$ \; \square \; $與斜步、多步較$ \; \square \; $相乘，
得$ \; \square \; $為同數，與左相消得$ \; \square \; $。
開平方得三百二十步，即大勾也。合問。}

\vspace{0.5cm}
\begin{center}
\longline
\end{center}

\wen{一百四十二}
\owen{%
甲乙二人共立於艮隅，乙南行過城門而立，甲東行望乙與城參相直而止。
丙丁二人共立於坤隅，丁向東行過城門而立，丙向南行望丁及甲乙悉與城俱相直。
丙復就甲斜行六百八十步與甲相會。
乙、丁又雲：「吾二人直行共得三百四十二步。」問答同前。}
\dayue{城徑二百四十步。}

\fayue{%
二雲數相乘，倍之為實；倍斜行於上，以二雲數相減加上位為從；一步常法。
開平方，得城徑。}

\tsaoyue{%
別得斜步即大弦也。其共步則一徑一虛弦共也，
其二數相並為一大和一虛弦共數也。
立天元為徑，減於共步得$ \; \square \; $為虛弦也。
以虛弦復減於天元，得$ \; \square \; $為虛和，以斜步乘之，得$ \; \square \; $（寄左）。
乃以天元加斜步，得$ \; \square \; $為大和，以虛弦乘之得$ \; \square \; $為同數，
與左相消得$ \; \square \; $。開平方得二百四十步，即城徑也。合問。}

\vspace{0.5cm}
\begin{center}
\longline
\end{center}

\wen{一百四十三}
\owen{%
甲從北門向東直行，庚從西門穿城東行，丙從西門向南直行，壬從北門穿城南行。
四人遙相望，悉與城參相直。
只雲甲丙相望處六百八十步，庚壬穿城共行了六百三十一步。問答同前。}
\dayue{城徑二百四十步。}

\fayue{%
共步自之得數。以共步減斜，餘自乘以減上為實。
二之斜步加入共步減斜，餘數為從。一步常法。得城徑。}

\tsaoyue{%
共行步為一徑與皇和共也，又為大和皇弦差也。甲丙相望即大弦也。
以共步減大弦，餘$ \; \square \; $為皇極弦上減一徑也。
立天元一為圓徑，減於共步，得$ \; \square \; $為皇極和也，以自之得$ \; \square \; $於上。
弦內減共步餘$ \; \square \; $，又以天元加之為皇弦，以自之得$ \; \square \; $，
減上位，餘得$ \; \square \; $為兩個皇直積（寄左）。
乃以天元乘皇弦得下式$ \; \square \; $為同數，與左相消得$ \; \square \; $。
平方而一，得二百四十步，即城徑也。合問。}

\newpage
\subsection{\tradfont 細草卷九下\quad 大和八問}

\begin{center}
\longline
\end{center}

\wen{一百四十四}
\owen{%
庚從西門穿城東行二百五十六步而立，壬從北門穿城南行三百七十五步而立。
又有甲丙二人俱在乾隅，甲向東行，丙向南行，各不知步數而立。
四人遙相望，只雲甲丙共行了九百二十步。問答同前。}
\dayue{城徑二百四十步。}

\fayue{%
庚東行冪、壬南行冪相並於上，並庚壬步而倍之，內減大和，
餘復減於庚壬共，得數以自乘減上位為平實。
並庚壬步為益從，半步為隅法。得城徑。}

\tsaoyue{%
立天元一為圓徑，以半之副置二位。
上以減於庚東行，得下$ \; \square \; $為平弦也；下以減於壬南行，得$ \; \square \; $為高弦也。
二弦相並得$ \; \square \; $為皇弦、虛弦共也。
倍此數得$ \; \square \; $為大弦、虛弦共也。
以大弦、虛弦共減於大和，餘$ \; \square \; $為虛勾虛股共也。
天元內減虛勾、虛股共，餘$ \; \square \; $即虛弦也。
復置皇弦虛弦共內減虛弦，餘$ \; \square \; $即皇極弦也，
以自之得$ \; \square \; $（寄左）。
然後以平弦自之，得下式$ \; \square \; $為勾冪也，又以高弦自之，
得$ \; \square \; $為股冪也。
二冪相並得$ \; \square \; $為同數，與左相消得$ \; \square \; $。
平方而一，得二百四十步，即城徑也。合問。}

\vspace{0.5cm}
\begin{center}
\longline
\end{center}

\wen{一百四十五}
\owen{%
甲丙俱在西北隅，甲向東行不知步數而立，丙向南行望見甲，就甲斜行與甲相會。
甲直行、丙直行共九百二十步（甲步少於丙步）。
又：出東門南行有柳樹一株，出南門東行有槐樹一株。
戊己二人同在巽隅，戊就柳樹，己就槐樹，亦與甲、乙遙相望。
只雲己行少於戊行數與兩樹相距數相並，得一百四十四步，其二數相減餘六十步。問答同前。}
\dayue{城徑二百四十步，太虛勾四十八步。}

\fayue{%
二雲數相並而半之為虛弦，以乘大和九百二十步於上。
以一百四十四減大和，以虛較乘之，減上位為平實。
以一百四十四減大和，又二之於上，以二之虛較減上位為從。四虛隅。得太虛勾四十八。}

\tsaoyue{%
別得甲丙直行共即大和也，戊就柳樹步即虛股也，己就槐樹步即虛勾也。
其一百四十四步即二明勾，其六十步即二股也。

立天元一為虛勾，加明勾得$ \; \square \; $為半徑也，
倍之得$ \; \square \; $即城徑也（又為虛弦上三事和）。
二雲數相並而半之，得$ \; \square \; $即小弦也。
相減而半之，得$ \; \square \; $即小較也。
以天元加較得$ \; \square \; $即小股也，小勾股共得$ \; \square \; $即小和也。
以小三事減大和得$ \; \square \; $即大弦也。
乃先置小和以大弦乘之，得下式$ \; \square \; $（寄左）。
次以小弦乘大和，得$ \; \square \; $，與左相消得下式$ \; \square \; $。
開平方得四十八步，即虛勾也。加明勾又倍之得二百四十步，即城徑也。合問。}

\vspace{0.5cm}
\begin{center}
\longline
\end{center}

\wen{一百四十六}
\owen{%
甲從乾隅東行，乙從艮隅南行，丙從乾隅南行，丁從坤隅東行，四人遙相望見。
既而甲還至艮隅就乙，丙還至坤隅就丁。
甲丙直行共九百二十步，甲還就乙共二百三十步，丙還就丁共五百五十二步。問答同前。}
\dayue{城徑二百四十步，虛弦一百二步。}

\fayue{%
並就數以減直行共，復以所並就數乘之為實；
並就數減直行共，得數復加入直行共為法。得虛弦。}

\tsaoyue{%
別得甲丙直行共為大和也，甲還就乙步為小差勾股共也，丙還就丁步為大差勾股共也。
以大差勾股共減於大股，餘即虛勾也。
以小差勾股共減於大勾，餘即虛股也。
二數相並得$ \; \square \; $為大弦虛弦共也，
二數相減，餘$ \; \square \; $為通差及太虛勾股差共也。
又並二數而半之，得$ \; \square \; $為皇極弦虛弦共，又為皇極勾股共也。

立天元一為虛弦。
先以二共數減於大和，餘$ \; \square \; $為虛勾虛股和於上。
次以虛弦減於二共數，餘$ \; \square \; $為大弦，
以乘上位，得下$ \; \square \; $（寄左）。
然後以天元乘大和，得$ \; \square \; $元為同數，與左相消得$ \; \square \; $。
上法下實，得一百二步，即虛弦也。加入虛和得二百四十步，即城徑也。合問。}

\fayue{%
又法：並雲數減大和，復以雲數相減乘之為實；
並雲數減大和，得數復加入大和為法。如法得虛差四十二。}

\tsaoyue{%
立天元一為虛較。
先以並雲數減大和，餘$ \; \square \; $為虛和於上。
次以天元減於$ \; \square \; $得$ \; \square \; $為通差，以乘之得$ \; \square \; $（寄左）。
然後以天元乘大和為同數，與左相消得$ \; \square \; $。
上法下實得四十二步，即虛差也。
副置虛和為二位，上加虛差而半之，得九十即虛股也。
下減虛差而半之，得四十八即虛勾也。
勾冪$ \; \square \; $，股冪得$ \; \square \; $，相並得$ \; \square \; $。
開平方得一百二步，即虛弦也。加入虛和得二百四十步，即城徑也。合問。}

\vspace{0.5cm}
\begin{center}
\longline
\end{center}

\wen{一百四十七}
\owen{%
依前見大和，只雲股圓差上勾弦差二百一十六，勾圓差上股弦差二十步。問答同前。}
\dayue{城徑二百四十步，小差股一百五十步。}

\fayue{%
以雲數二十步減通和，復以二十步乘之於上。
以雲數二百一十六減九百步而半之，乘上位為立實。
三因二十步以減通和得八百六十，
以二百一十六及二十共得二百三十六，減通和而半之，得三百四十二。
二數相乘訖，內減二十之九百步。
又以三百四十二及二百一十六共得五百五十八，又二十之以減之為從方。
以二百三十六減通和，又以三之二十步減通和，相並於上。
以二之五百五十八內卻減二十步，餘以加上位為益廉。
四步常法。得小差股一百五十。}

\tsaoyue{%
別得小差上股弦差$ \; \square \; $加二股為大勾也，大差上勾弦差$ \; \square \; $加二勾為大股也。

立天元一為小差股，加$ \; \square \; $得$ \; \square \; $為小差弦也，
小差弦上又加天元得$ \; \square \; $為通勾，以減於和步得$ \; \square \; $為通股也。
通股內減$ \; \square \; $得$ \; \square \; $，半之得下式$ \; \square \; $即大差之勾也。
大差勾上又加$ \; \square \; $，得$ \; \square \; $為大差弦也。
再置通股以小差弦乘之，得$ \; \square \; $，以天元除之，得$ \; \square \; $為一個大弦也（泛寄）。
再置通勾以大差弦乘之，得$ \; \square \; $，合以大差勾除。
不除寄為母，便以為大弦（寄左）。
乃以大差勾乘泛寄，得$ \; \square \; $為同數，與左相消得$ \; \square \; $。
益積開立方，得一百五十步，為小差股也。合問。}

\vspace{0.5cm}
\begin{center}
\longline
\end{center}

\wen{一百四十八}
\owen{%
依前見大和，只雲高弦、平弦共得三百九十一步，高弦、平弦相較得一百一十九步。問答同前。}
\dayue{城徑二百四十步。}

\fayue{%
以較數冪減於共數冪，又半之為實，以共數減大和為益從，一常法。
開平方，得圓徑。}

\tsaoyue{%
別得高弦減於通股為邊股內減明股也，平弦減於通勾為邊勾內減明勾也。
其共數即大弦內減皇極弦，又為皇極勾股共也，其相較步即皇極差也。
二雲數相並得$ \; \square \; $，即黃廣弦也。
二雲數相減，餘即黃長弦也。
以共數減於大和，餘$ \; \square \; $為皇極弦、圓徑共。

立天元一為圓徑，以減於$ \; \square \; $為皇極弦也。
以共數自之得$ \; \square \; $於上。
以相較數自之得$ \; \square \; $，減上位，餘$ \; \square \; $，
又半之得$ \; \square \; $為兩段皇極積（寄左）。
乃以天元乘皇極弦，得$ \; \square \; $為同數，
與左相消得下$ \; \square \; $。開平方得二百四十步，即城徑也。合問。}

\vspace{0.5cm}
\begin{center}
\longline
\end{center}

\wen{一百四十九}
\owen{%
依前見大和，只雲大差弦四百零八步，小差弦一百七十步。問答同前。}
\dayue{城徑二百四十步。}

\fayue{%
以並雲數減大和，復以相並數乘大和，又倍之為平實。
三之通和於上，又以並雲數減大和，加上位為從。二步虛法。得圓徑。}

\tsaoyue{%
大差弦減和步，餘$ \; \square \; $為大勾、大差勾共也。
以小差弦減大和，餘$ \; \square \; $為大股、小差股共也。
雲數相並得即大弦內減虛弦也，雲數相減得$ \; \square \; $為虛弦平弦共也。
以相並數減於大和，餘$ \; \square \; $為大差勾、小差股共，又為圓徑、虛弦共也。

立天元一為圓徑，減於$ \; \square \; $得$ \; \square \; $為虛弦也。
反以減於圓徑得$ \; \square \; $為小和也。
以天元減大和得$ \; \square \; $為大弦，以乘小和，得$ \; \square \; $（寄左）。
乃再置虛弦以通和乘之，得$ \; \square \; $，與左相消得$ \; \square \; $。
開平方得二百四十步，即城徑也。合問。}

\vspace{0.5cm}
\begin{center}
\longline
\end{center}

\wen{一百五十}
\owen{%
依前見大和，只雲黃廣弦五百一十步，黃長弦二百七十二步。問答同前。}
\dayue{城徑二百四十步，虛弦一百二步。}

\fayue{%
雲數相並減大和，復以相並數乘之為實。
雲數相並減大和，得數復以加大和為法。得虛弦一百二。}

\tsaoyue{%
別得黃廣弦又為大差弦、虛弦共，又為邊股、股共也。
黃長弦又為小差弦、虛弦共，又為底勾、明勾共也。
以黃廣弦減於大股，餘$ \; \square \; $即虛股，以黃長弦減於大勾，餘$ \; \square \; $即虛勾。
故並數以減於大和，餘$ \; \square \; $為虛和也。
以虛和減徑即虛弦也。
二雲數相並得$ \; \square \; $為大弦、虛弦共也，雲數相減餘$ \; \square \; $為虛弦、平弦共。

立天元一為虛弦，以減於七百八十二，得$ \; \square \; $為大弦也，
以小和乘之得$ \; \square \; $（寄左）。
乃以天元虛弦乘大和，得$ \; \square \; $為同數，與左相消得$ \; \square \; $。
上法下實，得一百二步，即虛弦也。合問。}

\vspace{0.5cm}
\begin{center}
\longline
\end{center}

\wen{一百五十一}
\owen{%
依前見大和，只雲邊弦五百四十四步，底弦四百二十五步。問答同前。}
\dayue{城徑二百四十步，皇極弦二百八十九步。}

\fayue{雲數相減自之為實，以大和減並數為法。得皇極弦。}

\tsaoyue{%
別得以邊弦減大股，餘$ \; \square \; $為半徑內減平勾，又為平弦內減勾圓差也。
以大勾減於底弦，餘$ \; \square \; $為高股內少半徑，又為股圓差內少高弦也。
二雲數相並，得九百六十九為大弦、皇極弦共也。
二雲數相減，得$ \; \square \; $為皇極勾股差也。
並數內減通和，餘$ \; \square \; $為皇極弦內減圓徑也。

立天元一為皇極弦，以自之於上。
以一百一十九自之得$ \; \square \; $，減上位得$ \; \square \; $為二皇積（寄左）。
復置天元內減四十九得下式$ \; \square \; $為黃方。
復以天元乘之，得$ \; \square \; $，與左相消得$ \; \square \; $。
上法下實，得二百八十九步，即皇極弦也。內減四十九，餘即城徑也。合問。}

\vspace{0.5cm}
\begin{center}
\longline
\end{center}

\wen{一百五十二}
\owen{%
依前見大和，只雲通弦六百八十步，通勾三百二十步。問答同前。}
\dayue{城徑二百四十步，通股六百步。}

\fayue{%
以通勾減通弦，餘以自之為平實。以通勾為從，一步常法。開平方得城徑。
又法：通弦冪內減通勾冪為實，以倍通勾為從，一步常法。得通股。}

\tsaoyue{%
別得通弦即大弦也，通勾即大勾也，通股即大股也。
其大勾股形為全書之宗，餘十四勾股形皆由此而生也。
立天元一為城徑。置通勾以天元乘之，又以通弦除之，得$ \; \square \; $為半徑之勾也。
又以通股乘之，得$ \; \square \; $，以通弦除之得$ \; \square \; $為半徑之股也。
勾股相乘得$ \; \square \; $，以天元半徑除之得$ \; \square \; $為弦也（寄左）。
然後以通弦$ \; \square \; $與左相消得$ \; \square \; $。
開平方得二百四十步，即城徑也。合問。}

\vspace{0.5cm}
\begin{center}
\longline
\end{center}

\wen{一百五十三}
\owen{%
依前見大和，只雲皇極弦二百八十九步，皇極勾股差一百一十九步。問答同前。}
\dayue{城徑二百四十步，皇極勾一百三十六步，皇極股二百五十五步。}

\fayue{%
以較冪減弦冪，又半之為實，弦內減較為從，一步常法。得皇極勾。
又以弦冪內減較冪為實，弦加較為從，一步常法。得皇極股。
又以弦內減較餘四之為實，一步常法。得黃方（即城徑減皇極弦也）。}

\tsaoyue{%
識別得皇極弦內減圓徑即皇極黃方也，皇極勾加皇極股即皇極和也。
立天元一為皇極勾，加較得$ \; \square \; $為皇極股，又加天元得$ \; \square \; $為皇極和也。
以和自之得$ \; \square \; $，內減弦冪$ \; \square \; $，餘$ \; \square \; $為二積（寄左）。
乃以天元乘皇極股得$ \; \square \; $，又倍之為同數，與左相消得$ \; \square \; $。
開平方得一百三十六步，即皇極勾也。
加較得二百五十五步，即皇極股也。
以弦冪減和冪，又半之得$ \; \square \; $為積，以天元半徑除之得$ \; \square \; $為和。
和內減弦餘即黃方，亦為城徑減皇極弦之數也。合問。}

\newpage

% ========================================================================
\section*{跋}
\addcontentsline{toc}{section}{跋}

\begin{tradfont}
測圓海鏡一書，自序雲「積思惑悟」，蓋先生潛心二十餘年而後成之者也。
其術以天元一為宗，太極為本，立一未知數而萬端可推，實開中國代數學之先河。
全書一百七十問，無一問不備問、答、術、草四體，條理密察，邏輯謹嚴，
雖古希臘之幾何原本，其精密不過如此。

卷七至卷九，計三十四問，方程式之高有至五乘、六乘者。
其開方之術，或益積，或翻法，或倒積，或消息，變化萬端，而歸於一理。
先生又創立十五勾股形之說，以通勾股為母，明勾股、、極勾股、虛勾股為子，
母子相生，條理粲然。此實中國數學史上之偉大創造也。

後之學者，如元之朱世傑、郭守敬，明之唐順之、顧應祥，清之李善蘭、華蘅芳，
皆深受此書之影響。李善蘭譯西算時，猶以天元術與代數相發明，
謂「中國之天元，即西國之代數」，其說誠不我欺也。

嗚呼！七百餘年過去，算板之制已湮，天元之術猶存。
讀此書者，不獨見中算之精深，亦可感先賢之睿智。
謹識數語，以為後來之勸云。
\end{tradfont}

\vspace{0.5cm}
\begin{center}
\longline
\end{center}

\newpage

% ========================================================================
\section*{Colophon}
\addcontentsline{toc}{section}{Colophon}

\begin{center}
\longline
\end{center}

\textit{測圓海鏡} (Ceyuan Haijing / Sea Mirror of Circle Measurements) was written by
Li Ye (李冶, 1192--1279) and completed in 1248. It is the earliest extant work
that systematically develops the \textit{tian yuan shu} (天元術) algebraic method
for solving polynomial equations.

The text uses a sophisticated system of 15 right triangles (十五勾股形) formed by
lines through and around a circular city, with standardized names for all edges
and diagonals. Each problem presents a configuration and asks for the city's
diameter (城徑), typically 240 paces (步), with the radius being 120 paces.

The mathematical expressions denoted by ■ in the original text represent
intermediate polynomial expressions arranged on the counting board according
to the ``Tai Ji'' (太極) and ``Tian Yuan'' (天元) positions---the ancient Chinese
equivalent of polynomial notation.

\textbf{Volumes 7--9 contain 34 problems:}
\begin{itemize}[leftmargin=2em]
  \item \textbf{Volume 7} (卷七): 18 problems on ``front'' configurations (明前一十八問),
    covering problems 105--122. These feature higher-degree polynomial equations
    (degrees 3--4) with detailed working on the \textit{tian yuan} method.
  \item \textbf{Volume 8} (卷八): 16 problems on ``rear'' configurations (明後一十六問),
    covering problems 123--138. These introduce multiple \textit{tian yuan} unknowns
    and cross-triangle relationships with equations up to degree 4.
  \item \textbf{Volume 9} (卷九): 12 problems on great oblique (大斜四問, problems 139--142)
    and great sum (大和八問, problems 143--150) configurations. These contain the
    most complex equations in the work, reaching degrees 5--6.
\end{itemize}

\vspace{0.5cm}
\textbf{Source:} Transcribed from Chinese Wikisource
(https://zh.wikisource.org/wiki/測圓海鏡).

\textbf{Expanded edition} adds detailed \textit{ts'ao} (草) working and \textit{ta} (答)
answer sections, with full polynomial derivations and expanded identification
passages (識別得) for all 34 problems.

\textbf{Typesetting:} Xe\LaTeX\ with xeCJK and Noto Sans CJK SC.

\end{document}
